% P10-4a, 14.09.2026: Mechanismus, Quellenreichweite und v4-Nachweis angeglichen.
% M08, 11.09.2026: Anlass und Entwurfsfolge; E47-02 mit Grenzen erhalten.
\section{Einordnende Fortschreibung des Projektbestands}
\label{sec:einordnende-fortschreibung}

\emph{Zielbild und Designproblem.} Der persistente Zustand aus
Abschnitt~\ref{sec:projektzustand-kontrollierte-mutation} ermöglicht
den Bezug auf bereits erarbeitetes Projektwissen. Eine neue Aussage
kann dieses bestätigen, präzisieren, ersetzen oder ihm widersprechen.
Eine rein additive Verarbeitung würde dagegen bei jedem Meeting
weitere Elemente anlegen und könnte konkurrierende Fassungen
desselben Sachverhalts erzeugen.

\emph{Entwurfsentscheidung.} Ein agentischer Einordnungsschritt
interpretiert neue Informationen deshalb relativ zum vorhandenen
Bestand und schlägt eine typisierte Einordnung vor: Neuzugang,
Bestätigung, Änderung, Widerspruch oder Klärungsbedarf. Auch eine
Bestätigung kann zusätzliche Herkunftsinformation liefern; deren
durchgängige Anbindung bleibt eine Implementierungsgrenze
(Abschnitt~\ref{sec:core-projektfortschreibung}). Widersprüche
und unbeantwortete Fragen bleiben als eigene
Entscheidungselemente sichtbar, statt zu einem vermeintlich
eindeutigen Ergebnis geglättet zu werden. Die resultierenden
Änderungen unterliegen den bestehenden Autorisierungspfaden.

Ein isolierter Experimentlauf illustriert den Bedarf an expliziter
Ablösung (E47-02). Einer bestehenden Anforderung, Ausschluss-Listen
je Bewohner zu führen, widersprach ein Folge-Meeting mit einer
übergreifenden Geltung. Die Einordnung erzeugte einen
Entscheidungsgegenstand; seine am Entscheidungs-Gate eingereichte
Auflösung führte zu einer neuen Anforderung mit Ablösungsbeziehung.
Die alte Fassung blieb historisch erhalten. Betroffene Arbeitselemente
wurden zur Angleichung vorgeschlagen, deren fachliche Bearbeitung
aber nicht abgeschlossen. Nur das Entscheidungs-Gate war interaktiv,
vier weitere Gates liefen im Experimentmodus; die anschließende
Rücksetzung ist nur narrativ dokumentiert. Der Fall zeigt den
ausgeführten Weg unter diesen Bedingungen. Seine fachliche Bewertung
folgt in Abschnitt~\ref{sec:eval-fortschreibung}.

Auch früher abgelehnte Vorschläge werden als Kontext für spätere
Einordnungen erhalten. Sie liefern Hinweise bei wiederkehrenden
Themen, lösen aber keine automatische Ablehnung aus; die erneute
Entscheidung bleibt beim Menschen. Die Gestaltung gilt für
Anforderungs- und Architekturinformationen. Daraus ergibt sich die
siebte Designanforderung:

\textbf{A7:} Neue fachliche Informationen müssen relativ zum
autorisierten Projektzustand als Neuzugang, Bestätigung, Änderung,
Widerspruch oder Klärungsbedarf eingeordnet werden; daraus resultierende
Zustandsänderungen müssen über die bestehenden kontrollierten
Fortschreibungs- und Autorisierungspfade erfolgen.

\emph{Einordnung und Grenzen.} Die Einordnung bleibt eine offene
semantische Aufgabe und kann fehlerhaft sein. Die beschriebenen
Kontrollen begrenzen ihre Zustandswirkung, belegen aber keine
allgemeine Erkennungsleistung für Bestandsinhalte oder Widersprüche.
Abschnitt~\ref{sec:core-projektfortschreibung} erklärt die resultierende
Fortschreibungslogik; Kapitel~\ref{chap:evaluation} untersucht
ausgewählte historische Fälle. Der folgende Abschnitt erweitert den
Kreislauf um operative Darstellungen und externe Rückmeldungen.
